% !TeX root = ../../main.tex
% !TEX root = ../main.tex

\section{\bilingual{Lemmas on \texorpdfstring{$N$}{N}-BRW}{有关\texorpdfstring{$N$}{N}-分支随机游走的引理}}\label{section:lemmas-on-n-brw}

\bilingual{
    We use the notation $\mathcal C$ to represent the set of all finite counting measures on $\mathbb R$. For $\mu\in\mathcal C$, we write $M(\mu)$ for the total mass of $\mu$, i.e., the number of particles.
}{
    记 $\mathcal C$ 为实数轴上所有有限计数测度的集合.~ 对于 $\mu\in\mathcal C$,~ 记 $M(\mu)$ 为 $\mu$ 的总质量,~ 即其粒子数.
}

\bilingual{
    We also write $\mu_1\succcurlyeq\mu_2$ if $M(\mu_1)\le M(\mu_2)$ and the atoms of $\mu_1$ are greater than or equal to those of $\mu_2$ in the sense that if $\mu_1=\sum_{i=1}^{M(\mu_1)}\delta_{x_i}$ and $\mu_2=\sum_{i=1}^{M(\mu_2)}\delta_{y_i}$ with $x_1\le x_2\le\cdots\le x_{M(\mu_1)}$ and $y_1\le y_2\le\cdots\le y_{M(\mu_2)}$, then $x_i\ge y_i$ for all $i\in[\![1,M(\mu_1)]\!]$. In the degenerate case where $M(\mu_1)=0$, we have $\mu_1\succcurlyeq\mu_2$ for any $\mu_2$.
}{
    称$\mu_1$占优于$\mu_2$,~ 记作$\mu_1\succcurlyeq\mu_2$,~ 若 $\mu_1$的粒子数不多于$\mu_2$,~ 即$M(\mu_1)\le M(\mu_2)$,~ 且 $\mu_1$ 的各原子位置依次大于等于 $\mu_2$ 靠左侧的各原子位置,~ 即,~ 若记 $$\mu_1=\sum_{i=1}^{M(\mu_1)}\delta_{x_i},\quad \mu_2=\sum_{i=1}^{M(\mu_2)}\delta_{y_i},$$ 其中 $x_1\le x_2\le\cdots\le x_{M(\mu_1)}$ 且 $y_1\le y_2\le\cdots\le y_{M(\mu_2)}$,~ 则 $x_i\ge y_i$ 对所有 $i\in[\![1,M(\mu_1)]\!]$ 成立.~ 在退化的情形 $M(\mu_1)=0$,~ 对任何 $\mu_2$,~ 都有 $\mu_1\succcurlyeq\mu_2$.
}

\bilingual{
    For all $N\ge 1$, let $\mathcal C(N)$ be the set of finite counting measures on $\mathbb R$ with total mass equal to $N$. We identify $\mathcal C(N)$ with the family of sets of real numbers with $N$ elements, and we endow $\mathcal C(N)$ with the partial order $\succcurlyeq$ defined above. For example, the population of particles $\mathcal Z^N_n$ at time $n$ in the $N$-BRW $\mathcal V^N$ can be identified with an element of $\mathcal C(N)$, and we write $\mathcal Z^N_n\in\mathcal C(N)$ for all $n\ge 0$. Recall that we may also write $\mathcal V^N_n$ in lieu of $\mathcal Z^N_n$ to represent the population of particles at time $n$ in $\mathcal V^N$.
}{
    对任一 $N\ge 1$,~ 记 $\mathcal C(N)$ 为实数轴上总质量为 $N$ 的有限计数测度全体组成的集合.~ 并且,~ 对于$\mathcal C(N)$中的某一测度,~ 它可被视作为元素个数为 $N$ 的实数集合.~ 例如,~ 在 $N$-分支随机游走 $\mathcal V^N$ 的术语中,~ 时刻 $n$ 的粒子全体 $\mathcal Z^N_n$ 可视作为 $\mathcal C(N)$ 中的一个测度,~ 此时可记 $\mathcal Z^N_n\in\mathcal C(N)$.~ 回顾上文,~ 也可以用 $\mathcal V^N_n$ 来表示时刻 $n$ 的粒子全体 $\mathcal Z^N_n$.~ 另外,~ 在 $\mathcal C(N)$ 上也可类似定义上述占优关系 $\succcurlyeq$.
}

\bilingual{
    We also denote by $\max\mu$ and $\min\mu$ the maximum and minimum of the support of $\mu\in\mathcal C$, respectively. For example, the leftmost position $\mathcal M^N_n$ at time $n$ in the $N$-BRW $\mathcal V^N$ can be identified with $\min\mathcal Z^N_n$.
}{
    对 $\mu\in\mathcal C$,~ 记 $\max\mu$ 和 $\min\mu$ 分别为 $\mu$ 的支撑集的最大值和最小值.~ 例如,~ 在 $N$-分支随机游走 $\mathcal V^N$ 的术语中,~ 时刻 $n$ 的最左侧位置 $\mathcal M^N_n$ 可视作为 $\min\mathcal Z^N_n$.
}
